\documentclass[11pt,a4paper]{article}
\usepackage[utf8]{inputenc}
\usepackage[T1]{fontenc}
\usepackage[spanish]{babel}
\usepackage{geometry}
\geometry{margin=2.5cm}
\usepackage{setspace}
\setstretch{1.08}
\usepackage{hyperref}
\usepackage{longtable}
\usepackage{array}
\usepackage{booktabs}
\usepackage{listings}
\usepackage{xcolor}
\usepackage{enumitem}
\hypersetup{colorlinks=true,linkcolor=blue,urlcolor=blue}

\lstdefinestyle{cmd}{
  basicstyle=\ttfamily\small,
  frame=single,
  breaklines=true,
  showstringspaces=false
}

\title{Meso2D: Tutorial de Usuario}
\author{Guia practica para usuarios nuevos}
\date{\today}

\begin{document}
\maketitle
\tableofcontents
\newpage

\section{Objetivo de este tutorial}
Este documento explica como utilizar Meso2D para generar mesoestructuras 2D de materiales con:
\begin{itemize}[leftmargin=1.5em]
  \item aridos (circulares, elipticos o poligonales),
  \item poros,
  \item puntos reactivos,
  \item y exportacion de resultados para postproceso.
\end{itemize}

El enfoque es practico: seguir pasos claros para ejecutar simulaciones reproducibles y exportar datos correctos.

\section{Que hace el programa}
Meso2D permite construir una probeta 2D y poblarla con entidades geometricas de forma controlada por:
\begin{itemize}[leftmargin=1.5em]
  \item dosificacion granulometrica,
  \item fraccion de area de aridos,
  \item fraccion de area de poros,
  \item fraccion de area de puntos reactivos,
  \item y semilla aleatoria (seed).
\end{itemize}

El flujo interno esta orquestado por un controlador de simulacion que selecciona el worker segun el modo:
\begin{itemize}[leftmargin=1.5em]
  \item \textbf{circulos},
  \item \textbf{ellipses},
  \item \textbf{poligonos}.
\end{itemize}

\section{Requisitos previos}
\subsection{Dependencias}
Asegurate de tener instalados al menos:
\begin{itemize}[leftmargin=1.5em]
  \item Python 3.x
  \item PyQt5
  \item numpy
  \item matplotlib
  \item ezdxf
  \item shapely
  \item opensimplex
\end{itemize}

\subsection{Ejecucion}
Desde la carpeta del proyecto, la entrada principal recomendada es:
\begin{lstlisting}[style=cmd]
python Meso2d.py
\end{lstlisting}

Alias legacy de compatibilidad (siguen funcionando):
\begin{lstlisting}[style=cmd]
python Meso2D_refactor.py
python Meso2D_refactor_EN.py
\end{lstlisting}

\section{Flujo de trabajo rapido (5 pasos)}
\begin{enumerate}[leftmargin=1.7em]
  \item Elegir modo geometrico (circulos, elipses o poligonos).
  \item Introducir dimensiones de probeta y parametros de dosificacion.
  \item Activar opcionalmente poros y/o puntos reactivos.
  \item Ejecutar la simulacion y revisar la visualizacion.
  \item Exportar estructura e imagen final.
\end{enumerate}

\section{Pantalla inicial: seleccion de modo}
Al iniciar, aparece una ventana de seleccion. Ese modo determina:
\begin{itemize}[leftmargin=1.5em]
  \item el algoritmo de colocacion de aridos,
  \item la representacion geometrica de salida,
  \item y el tipo de comprobaciones de colision.
\end{itemize}

\textbf{Recomendacion:}
\begin{itemize}[leftmargin=1.5em]
  \item Empieza con \textit{circulos} para pruebas rapidas.
  \item Usa \textit{ellipses} o \textit{poligonos} para estudios mas realistas.
\end{itemize}

\section{Configuracion de un proyecto}
\subsection{Parametros principales}
\begin{longtable}{p{0.28\textwidth} p{0.66\textwidth}}
\toprule
\textbf{Parametro} & \textbf{Significado y efecto} \\
\midrule
x, y & Dimensiones de la probeta (ancho y alto). Determinan el area total A = x*y. \\
Pagg & Fraccion de area objetivo ocupada por aridos. \\
Tabla sieve\_size / tpp & Define la curva granulometrica. El programa reparte area por fracciones usando estas columnas. \\
dporo\_min, dporo\_max & Diametros minimo y maximo para poros. \\
Pporos & Fraccion de area objetivo para poros. \\
dpto\_min\_aggregates, dpto\_max\_aggregates & Rango de diametros de puntos reactivos sobre aridos. \\
Ppto\_react\_aggregates & Fraccion de area de puntos reactivos en aridos. \\
dpto\_min\_pasta, dpto\_max\_pasta & Rango de diametros de puntos reactivos en pasta. \\
Ppto\_react\_pasta & Fraccion de area de puntos reactivos en pasta. \\
seed & Semilla aleatoria para reproducibilidad exacta. \\
\bottomrule
\end{longtable}

\subsection{Tabla de dosificacion}
\begin{itemize}[leftmargin=1.5em]
  \item Usa \textbf{Add} para insertar filas.
  \item Completa \textbf{sieve\_size} (tamano de tamiz) y \textbf{tpp} (porcentaje pasante acumulado).
  \item Mantener coherencia fisica evita fallos de convergencia.
\end{itemize}

\textbf{Consejo practico:}
si una simulacion tarda demasiado o no logra colocar entidades, reduce densidades objetivo (Pagg/Pporos/Ppto\_react\_*) o amplía la probeta.

\section{Opciones de simulacion}
\subsection{Poros}
Al activar la opcion de poros:
\begin{itemize}[leftmargin=1.5em]
  \item se calcula una lista de radios hasta consumir el area objetivo,
  \item y luego se colocan evitando colisiones con aridos y otros poros.
\end{itemize}

\subsection{Puntos reactivos}
Los puntos se pueden activar en dos canales:
\begin{itemize}[leftmargin=1.5em]
  \item \textbf{sobre aridos},
  \item \textbf{en pasta}.
\end{itemize}

Internamente:
\begin{itemize}[leftmargin=1.5em]
  \item para aridos se usa muestreo por \textbf{clusters},
  \item para pasta se usa \textbf{ruido Simplex} para patrones organicos no uniformes.
\end{itemize}

\section{Ejecucion de la simulacion}
\subsection{Antes de pulsar Run}
Checklist recomendado:
\begin{enumerate}[leftmargin=1.7em]
  \item Modo seleccionado correctamente.
  \item x, y con valores validos.
  \item Tabla de dosificacion completa.
  \item Pagg definido.
  \item Si hay poros/puntos, todos sus campos necesarios rellenados.
  \item Seed fijada si necesitas reproducibilidad.
\end{enumerate}

\subsection{Durante la ejecucion}
\begin{itemize}[leftmargin=1.5em]
  \item La barra de progreso avanza por etapas.
  \item El panel de log muestra estado y mensajes del worker.
  \item Puedes usar \textbf{Stop} para cancelar una corrida.
\end{itemize}

\subsection{Despues de la ejecucion}
Al terminar:
\begin{itemize}[leftmargin=1.5em]
  \item se habilita la visualizacion de estructura,
  \item se habilita la exportacion,
  \item y queda registrado el tiempo de simulacion.
\end{itemize}

\section{Visualizacion y guardado de imagen}
\begin{itemize}[leftmargin=1.5em]
  \item \textbf{Ver estructura}: abre el visor.
  \item \textbf{Save Image}: guarda en PNG/JPEG/TIFF/BMP.
\end{itemize}

\textbf{Nota:} el guardado esta pensado para conservar la estructura generada completa.

\section{Exportacion de estructura}
\subsection{Formatos disponibles}
Desde \textbf{Export Structure} puedes guardar en:
\begin{itemize}[leftmargin=1.5em]
  \item DXF
  \item SVG
  \item JSON
  \item Gmsh GEO
\end{itemize}

\subsection{Recomendaciones}
\begin{itemize}[leftmargin=1.5em]
  \item Exporta inmediatamente despues de una corrida correcta.
  \item Si cambias parametros, vuelve a ejecutar antes de exportar.
  \item Mantén nombres de archivo con fecha/seed para trazabilidad.
\end{itemize}

\section{Proyectos: nuevo, abrir, guardar, ejemplo}
\begin{itemize}[leftmargin=1.5em]
  \item \textbf{New Project}: limpia estado actual.
  \item \textbf{Open Project}: carga archivo de proyecto y rellena la UI.
  \item \textbf{Save Project}: guarda parametros actuales.
  \item \textbf{Load Example}: carga un caso de ejemplo para pruebas rapidas.
\end{itemize}

\section{Reproducibilidad y trazabilidad}
Para repetir exactamente una estructura:
\begin{enumerate}[leftmargin=1.7em]
  \item guarda el proyecto,
  \item anota la seed,
  \item mantén el mismo modo,
  \item usa la misma version del codigo.
\end{enumerate}

\section{Solucion de problemas comunes}
\subsection{La simulacion no arranca}
Causas habituales:
\begin{itemize}[leftmargin=1.5em]
  \item falta algun campo obligatorio,
  \item tabla de dosificacion incompleta,
  \item valores no numericos.
\end{itemize}

\subsection{No se colocan suficientes entidades}
\begin{itemize}[leftmargin=1.5em]
  \item baja Pagg, Pporos o Ppto\_react\_*,
  \item amplia dominio (x, y),
  \item reduce diametros minimos,
  \item simplifica inicialmente usando modo circulos.
\end{itemize}

\subsection{Exportacion falla}
\begin{itemize}[leftmargin=1.5em]
  \item verifica que exista una estructura recien generada,
  \item revisa permisos de escritura en carpeta destino,
  \item prueba formato DXF como control inicial.
\end{itemize}

\section{Buenas practicas para usuario final}
\begin{enumerate}[leftmargin=1.7em]
  \item Trabaja con versionado de proyectos (v1, v2, v3...) y seed registrada.
  \item Cambia un bloque de parametros cada vez para entender su impacto.
  \item Guarda imagen y estructura en cada corrida util.
  \item Documenta en un cuaderno: modo, seed, Pagg, Pporos, ratios reactivos.
  \item Para lotes de pruebas, empieza en circulos y pasa a elipses/poligonos al final.
\end{enumerate}

\section{Mini guia de inicio rapido}
\begin{enumerate}[leftmargin=1.7em]
  \item Ejecuta \texttt{python Meso2D\_refactor.py}.
  \item Elige modo \textit{circulos}.
  \item Define x=100, y=100 (ejemplo), Pagg y tabla de dosificacion.
  \item Activa poros con valores moderados.
  \item Pulsa \textbf{Run}.
  \item Visualiza resultado y exporta a DXF + PNG.
\end{enumerate}

\section{Notas para futuros usuarios avanzados}
\begin{itemize}[leftmargin=1.5em]
  \item El modo poligonal y eliptico usa comprobaciones geometricas mas exigentes.
  \item Los puntos reactivos en pasta dependen del umbral y escala del campo Simplex.
  \item En dominios muy saturados, el limite de intentos evita bucles infinitos y puede descartar candidatos.
\end{itemize}

\section{Resumen final}
Meso2D combina dosificacion, muestreo estocastico y restricciones geometricas para crear mesoestructuras 2D exportables.

Si sigues el flujo:
\begin{enumerate}[leftmargin=1.7em]
  \item configurar parametros,
  \item ejecutar,
  \item validar visualmente,
  \item exportar,
  \item registrar seed y proyecto,
\end{enumerate}
obtendras resultados reproducibles y comparables para analisis posteriores.

\end{document}
